\begin{lemma}[Existence of the canonical hierarchy base]
Let $T$ be any positive-rung extended Tashkinov tree, let $n$ be its last
rung, let $T_n$ be its terminal ladder tree, and let $\varphi_n$ be its
generating coloring.  Then $T$ admits a
\noderef{HierarchyBase} under $\varphi_n$.  No maximum-property, existing
good hierarchy, or lower bound on $\lvert\overline{\varphi_n}(T_n)\rvert$
is required.
\end{lemma}
\begin{proof}
By \noderef{ExtendedTashkinovTree}, the certified series has at least one
noninitial iteration and $\varphi_n$ is exactly the coloring recorded at
its last rung.  The certification field of \noderef{TashkinovSeries}
supplies a \noderef{TashkinovSeriesStep} certificate at every index; inspect
the one at the final index.  Its initial constructor is impossible:
that constructor certifies index $0$, whereas the final index is the
positive rung $n$.  We treat the remaining three constructors.

Suppose first that the final constructor is SE or RE.  By
\noderef{TreeSequence}, $T_n$ has positive length, so position $1$ is a
valid positive vertex position of $T_n$.  Take both the core and the
otherwise unused closure field to be $T_n$, and take support position $1$.
The terminal missing-color equality required by \noderef{HierarchyBase}
is reflexive because the base coloring is $\varphi_n$.  The constructor's
exact kind equation says respectively $\Theta_n=\mathrm{SE}$ or
$\Theta_n=\mathrm{RE}$, so the parallel implication in the definition is
vacuous.  If $\Theta_n\ne\mathrm{PE}$, the chosen core is exactly $T_n$,
which proves the nonparallel clause.  These data therefore form the
required base in both nonparallel cases.

It remains to consider a final PE constructor.  Let
$(\pi,v,f,u,\delta)$ be its recorded
\noderef{MaximumDefectiveChoice} on $T_n$.  We first prove that the
selected supporting vertex $v$ occurs at a positive position of $T_n$.
By \noderef{DefectiveEdge}, the exact selected edge identity $f$ belongs
to $\partial_{\pi,\alpha}(V(T_n))$ for some defective color $\alpha$.
By \noderef{DefectiveColor} and \noderef{BoundaryColorEdges}, that finite
set contains at least two distinct exact edge identities.  Choose
\[
              g\in\partial_{\pi,\alpha}(V(T_n)),\qquad g\ne f.  \tag{1}
\]
The definition of \noderef{BoundaryEdges} orients the ends of $g$ as
\[
       \operatorname{ends}(g)=\{w,z\},\qquad
       w\in V(T_n),\quad z\notin V(T_n).                         \tag{2}
\]
Because $g$ has the same defective color $\alpha$, it is itself a
defective edge, and (2) makes it incident with $w$.  Hence $w$ is a
defective vertex by \noderef{DefectiveVertex}.

Apply the maximality clause of \noderef{MaximumDefectiveChoice} to the
same admissible coloring $\pi$ and the defective vertex $w$.  It gives
either $w=v$ or $w\prec_{T_n}v$.  The first alternative is impossible.
Indeed, (1) gives two distinct exact edges $f$ and $g$; both have color
$\alpha$, and if $w=v$ then both are incident with $v$.  They are thus
adjacent in the sense of \noderef{AdjacentEdges}, contradicting the
properness of $\pi$ in \noderef{IsProperEdgeColoring}.  Consequently
\[
                              w\prec_{T_n}v.                     \tag{3}
\]
By \noderef{TreeVertexOrder}, (3) supplies vertex positions $r<s$ with
$T_n(r)=w$ and $T_n(s)=v$.  In particular $s\geq1$.  Take $s$ as the
support position.  The PE constructor's exact current-vertex equation
identifies the recorded chosen vertex with $v=T_n(s)$, as required by the
parallel clause of \noderef{HierarchyBase}.

Form the exact ordered prefix $T_n(v)$ ending at position $s$ using
\noderef{TreeSequencePrefix}; its definition applies because
$1\leq s\leq\lvert T_n\rvert$.  By
\noderef{TashkinovClosureExists}, choose a closure $R_n$ of this prefix
under the generating coloring $\varphi_n$.  Apply
\noderef{OrderedTreeUnionExists} to $T_n$ and $R_n$.  It yields a
tree-sequence $U$ and an exact ordered-union certificate
\[
             \noderef{OrderedTreeUnion}(T_n,R_n,U).              \tag{4}
\]
Set the hierarchy-base closure to $R_n$ and its core to $U$.  The support
position is positive by (3), the PE chosen-vertex equation was verified
above, the required prefix closure is the chosen closure, and (4) is the
required ordered union.  Again
$\overline{\varphi_n}(T_n)=\overline{\varphi_n}(T_n)$ is reflexive.  The
nonparallel implication is vacuous because the final constructor records
$\Theta_n=\mathrm{PE}$.  Thus all fields of \noderef{HierarchyBase} are
satisfied in the PE case as well, completing the construction.
\end{proof}
